% PLAN (20-25 pages). The paper text on data and methods is input in full; this chapter adds what the paper has no room for.
%   3.1 Data: paper/sections/data.tex (input below, brings its own \section). Add per-set preparation: file formats, preprocessing, orientation
%       checks, exclusions.
%   3.2 Methods: paper/sections/methods.tex (input below, brings its own \section; compact version of CT decomposition, projection, detector
%       physics, suppressor, v1 supervision, audit definition, real-image protocol, statistics). When the thesis text outgrows it, split methods.tex
%       into one file per subsection (in the paper folder) and input them here separately; do not fork the text.
%   3.3 Thesis-only expansion (below): segmentation recipe in detail, rendering geometry and frame, attenuation tables (ICRU-44
%       compositions, NIST tables with XCOM berger2010xcom), spectra (SpekPy poludniowski2021spekpy), beam hardening and the additivity gap, tone curves and
%       exactness of X' = S' + B', detector effects (blur, noise, scatter), renderer tests (tests/test_renderer.py, tests/test_compose.py),
%       orientation check, LUNA16 versus DICOM voxel check, usable-CT rules and the padded-row decision (DEVIATIONS.md, 2026-10-27).
%   Every statement in Section sec:dp-details is read from the code (s4/, tools/, kernels/, tests/); no result of the study is typed here.
This chapter first reproduces the data and methods text of the paper and then adds the details that the paper has no room for.
Section~\ref{sec:dp-details} extends the paper text below without repeating it: each of its subsections gives the implementation behind one
step of Sections~\ref{sec:data} and~\ref{sec:methods} as it stands in the code.

\section{Data}
\label{sec:data}
% PLAN (about 1 page and one table). Physics and rendering are in methods.tex; this section lists what data are used for what, under which licence,
% and which sets are locked.
%   3.1 Training sources: LIDC-IDRI (v1 renders; lesion library for the twin pastes; synthetic test CTs) and the public Qure ct2xr-projections
%       (v2 source; exactly additive components, see Section method_v2). Licence consequence: weights trained on the projections inherit CC BY-NC-SA 4.0.
%   3.2 Real evaluation data: SZCH-X-Rays (paired DES), VinDr-RibCXR (rib masks, licensed, local only), JSRT, Node21, TBX11K, VinBigData, NIH ChestX-ray14.
%   3.3 Data discipline: development versus locked sets, split rules, what was seen before the freeze (DEVIATIONS.md 2026-10-27, SAP_v2.md).
%   Table~data: dataset, role, licence or access. Licences are copied only from pages that were read; everything else stays a \todo.
% Counts below are copied from SAP_v2.md, methods.tex and DEVIATIONS.md; regenerate them from results/ when the tables are final.
Table~\ref{tab:data} lists every data set with its role in this study and its licence or access condition.

\subsection{Training sources}
\paragraph{LIDC-IDRI.} The LIDC-IDRI collection \citep{armato2011lidc,armato2015lidctcia,clark2013tcia}, with some series read from the LUNA16 copy
\citep{setio2017luna16}, supplies the v1 renders, the lesion library used to paste lesions into training images, and the synthetic test CTs. The
selection, the usable-CT rules and the patient-level split are given in Section~\ref{sec:methods}, together with the detector physics (attenuation
tables from the NIST X-ray mass attenuation coefficient database \citep{hubbell2004xraymass}, tube spectra from SpekPy \citep{poludniowski2021spekpy}).

\paragraph{Qure.ai CT projections.} The v2 models are trained on the public Qure.ai \path{ct2xr-projections} data \citep{angaitkar2026anatomy}, which
hold, for each CT volume, a composite image and translated bone, soft-tissue and lung projections. The composite is an exact linear combination of
the components up to rounding, so the bone image of every training sample is known exactly (Section~\ref{sec:method_v2}).
The release holds \num{21887} projections, one per CT volume; the volume names are those of CT-RATE \citep{hamamci2026ctrate}. We split them by
patient (\num{14374} patients) into \num{21051} training, 406 validation and 430 test projections (\path{tools/qp_prep.py}).

\subsection{Real radiographs}
\paragraph{SZCH-X-Rays.} Paired radiographs and dual-energy soft-tissue images from one partner hospital \citep{sun2026bsldm,sun2026debonedit}. We use the
75-pair test subset; its 75 images are split into a development part (24 images) and a held-out part (51 images) by a hash of the file name
(\path{tools/szch_split.py}). The 75 images are the published test split of the 741-pair version of the data set (592/74/75), acquired with dual-exposure DES on a GE
Discovery XR656 \citep{sun2026debonedit}; an earlier version of 818 pairs was used for BS-LDM \citep{sun2026bsldm}, and the papers do not
explain the difference.

\paragraph{VinDr-RibCXR.} 245 radiographs with per-rib masks \citep{nguyen2021ribcxr}. The data are licensed and are processed only on the licence
holder's own computer and private cloud account; only metrics leave them.

\paragraph{Nodule and detection sets.} JSRT \citep{shiraishi2000jsrt} (154 nodules), Node21 \citep{sogancioglu2024node21} (1\,476 nodule boxes),
TBX11K \citep{liu2020tbx11k} and VinBigData (the Kaggle competition data; related to VinDr-CXR \citep{nguyen2022vindrcxr}) are used for real-lesion contrast
and for downstream detection; NIH ChestX-ray14 \citep{wang2017chestxray} is used for downstream classification with its official lists.
NIH uses the official test list (\num{25596} images) and a patient-level split of the remaining images into \num{77890} training and \num{8634}
validation images. The detection sets are split by a fixed hash of the image identifier: Node21 \num{3440}/\num{512}/\num{930}, TBX11K
\num{6200}/\num{688}/\num{2088} (its official validation set serves as test set) and VinBigData \num{10433}/\num{1469}/\num{2943} images for
training, validation and test. VinBigData boxes are kept where at least two radiologists marked the same finding; images with only single-reader
findings are left out. No image of any set was unreadable. VinBigData is the Kaggle competition built on the VinDr-CXR images \citep{vinbigdata2020kaggle}; its rules allow non-commercial research use.

\paragraph{BDCXR-3257 (external validation).} 3,257 paediatric chest radiographs (JPEG) from the Radiology and Imaging Department of
Dr.~M~R~Khan Shishu Hospital \& Institute of Child Health, Dhaka, collected in routine care and anonymised, with one label each: normal (880),
bacterial pneumonia (1,477) or viral pneumonia (900) \citep{durjoy2026bdcxr}, a task studied before on other paediatric sets \citep{kermany2018identifying}. Most children are under five years old
(1,132 images under one year). The data were used under the ethical approval REC-FSIT-2025/No:~11928 (Daffodil International University) and
the data use agreement between the hospital and the authors' research group; no radiograph of this set is shown in this paper. There are no
patient identifiers, so near-duplicate images are grouped by their correlation at $64$~px ($r\ge0.93$; 48 pairs, the rest single images) and
a group never spans two folds.

\subsection{Development and locked sets}
SAP.md fixed the locked sets before any locked-set evaluation. After the v1 ablation the SZCH pairs were looked at (DEVIATIONS.md, 2026-10-27), so SZCH
became a development set for decisions about the v2 recipe, and SAP\_v2.md was written and frozen before any held-out or locked evaluation of a v2
model. Decisions used the 24 development images, the synthetic validation data and the real-image audit on all 75 SZCH images; the audit on SZCH is
therefore development evidence only, and the confirmatory audit runs on VinDr-RibCXR (not opened by any model run before the plan was written) and
JSRT. Paired fidelity claims use the 51 held-out SZCH pairs. Before the freeze, screening tools printed held-out PSNR and SSIM next to the
development numbers for v1-family and screening runs; no rule, table or decision used them, and the tool was changed to print development
numbers only.

\begin{table}[htbp]
\centering\small
\renewcommand{\arraystretch}{1.15}
\caption{Data sets, their role in this study and the licence or access condition, as stated on the providers' pages (read on
3 November 2026).}
\label{tab:data}
\begin{tabular}{@{}p{0.24\linewidth}p{0.35\linewidth}p{0.33\linewidth}@{}}
\toprule
Data set & Role here & Licence or access \\
\midrule
LIDC-IDRI (TCIA) & v1 renders, lesion library, synthetic test CTs & CC BY 3.0; attribution with the data DOI \\
Qure.ai ct2xr-projections & v2 training source & CC BY-NC-SA 4.0 \\
CT-RATE & CT source of the projections (not used directly) & CC BY-NC-SA 4.0, gated \\
SZCH-X-Rays & paired DES evaluation (24 development, 51 held-out) & CC BY-NC 4.0 \\
VinDr-RibCXR & locked audit and rib benchmark (245 images) & licensed; local use, metrics only \\
JSRT & audit and nodule contrast & free after registration; not public domain, no redistribution or commercial use \\
Node21 & nodule contrast, detection & CC BY-NC-ND 4.0 \\
TBX11K & detection & CC BY 4.0 (data) \\
VinBigData / VinDr-CXR & detection & competition rules: non-commercial research use \\
NIH ChestX-ray14 & classification & unrestricted use with attribution (NIH Clinical Center) \\
BDCXR-3257 & external validation (pneumonia classification, audit) & CC BY 4.0 (figshare); used under the providers' data use agreement \\
\bottomrule
\end{tabular}
\end{table}


\section{Methods}
\label{sec:methods}

\subsection{CT decomposition and projection}
\paragraph{Data.} We used the LIDC-IDRI collection (1\,018 chest CT series from 1\,010 patients) with the radiologists' nodule contours.
Series were read from DICOM or from the LUNA16 MetaImage copy (voxel-identical on all 12 series we compared, including nodule slice locations); 1\,000 series
could be rendered (18 were skipped because of non-axial orientation or irregular slice spacing). A render was \emph{usable} if the lungs were not cut by the
scan, at most 50 of the 512 frame rows lay beyond the scan, the lung volume was between 1.5 and 9~L and the bone volume exceeded 300~cm$^3$
(844 usable CTs). Patients were assigned to train/validation/test (70/10/20\%) by a hash of the patient identifier (588/98/158 CTs; 1\,057/172/298
consensus nodules) before any model was trained; the split file hash is recorded in \path{freeze.json}.

\paragraph{Components.} Each CT was split into bone, non-lung soft tissue and lung. The body is the largest component of $\mathrm{HU}>-300$
(slice-wise hole filling, then a $1{\times}9{\times}9$ opening to detach the couch). Bone follows the iterative-relaxation recipe of the reference method \citep[Algorithm~1]{angaitkar2026anatomy}, with a 10~mm support cube,
the first pass starting from the reference mask and a single closing after the ladder:
a reference mask at 300~HU, thresholds $\tau\in\{200,150,100\}$~HU admitted only where at least 5\% of a 10~mm neighbourhood was already bone, union with the
reference mask, a radius-2 closing to fill marrow, and removal of components smaller than 1\% of the largest. Lung is air ($<-320$~HU) inside the body in
components at least 20\% the size of the largest, closed and hole-filled slice-wise (so vessels and nodules belong to the lung component) and
without bone. Bone voxels were replaced in the soft-tissue component by a normalised convolution of the surrounding soft tissue (Gaussian, 4~mm, with
12 and 36~mm fall-backs). Density is water-equivalent, $\rho=\mathrm{clip}(1+\mathrm{HU}/1000,0,3)$~g\,cm$^{-3}$.

\paragraph{Projection.} Each component was integrated along the anteroposterior axis (parallel beam), giving an areal density $a_c$ (g\,cm$^{-2}$).
The frame is a body crop: width 1.08 times the body width, height the lung extent plus 60~mm clipped to 0.75--1.2 of the width (short scans cannot
reach 1.2), area-resized to $512\times512$ and later squashed to a square exactly as real radiographs are. Padded rows are set to air.
The AP view is the horizontal mirror of the PA view.

\paragraph{Detector physics.} A pixel value is
\begin{equation}
P_c=-\ln\frac{\sum_E \phi(E)\,E\,\exp\!\big(-(\mu/\rho)_c(E)\,a_c\big)}{\sum_E\phi(E)\,E},
\label{eq:physics}
\end{equation}
with ICRU-44 \citep{icru1989report44} cortical bone, soft tissue and lung attenuation (table~4 of NIST SRD 126; \citealp{hubbell2004xraymass}) and W-anode spectra (SpekPy, 12$^\circ$, 2.5~mm Al) at
70--140~kVp (training draws 80--130~kVp, the synthetic evaluation uses 100~kVp). Components are therefore stored as areal densities and the
spectrum is applied at training time, so one render serves every tube voltage.
The sum of separately projected components differs from the projection of the whole body because of beam hardening; we report this
additivity gap (6.1\% at 100~kVp for 20~g\,cm$^{-2}$ of soft tissue and 3~g\,cm$^{-2}$ of bone).

\subsection{Training images with exact components}
A training sample is composed on the GPU from the stored areal densities in this order: lesion insertion (below), one random affine transform shared by all maps
(rotation $\pm15^\circ$, scale 0.9--1.15, shift $\pm6\%$, mirror), the physics above, component weighting
$S=\mathrm{norm}(P_{\mathrm{nl}}+g\,P_{\mathrm{lung}})$, $g\sim U(1,3)$, $B=w\,P_{\mathrm{bone}}/p_{99}$, $w\sim U(0.2,0.5)$,
$X=(S+B)/1.5$, and a monotone tone curve $\gamma$. The tone curve is either a random gamma or histogram matching to the intensity quantiles of a randomly
drawn real radiograph (4\,000 NIH ChestX-ray14 training-list images). Defining $X'=\gamma(X)$, $S'=\gamma(S)$ and $B'=X'-S'$ keeps the
decomposition exact under any monotone mapping. Gaussian blur is linear and applied identically to $X'$ and $B'$; quantum noise and scatter
modify only $X'$.

\subsection{Suppressor}
The suppressor is a SwinV2-Tiny encoder (ImageNet-1k, window 8, $512^2$ input) with a nested (UNet++) decoder over the four Swin stages, CBAM attention, a plain
U-Net tail with two shallow convolutional skips to full resolution and a $1\times1$ head that outputs the bone image directly; the suppressed image is
$x-\hat b$. Relative to the reference this is a smaller backbone at half the resolution, forced by compute. Training used AdamW (lr $2\times10^{-4}$,
weight decay 0.01, 500-step warm-up, cosine decay), batch 8, mixed precision, gradient clipping at 1, EMA weights and 200\,000 samples per arm.
The loss is $0.5\,\mathrm{L1}+0.5\,\mathrm{MSE}+0.5\,(1-\mathrm{SSIM})+2.5\,(\mathrm{L1}+\mathrm{MSE})_{\mathrm{ROI}}$ with the ROI taken from the lung component
(top structures: $P_{\mathrm{lung}}/p_{99.9}>0.3$). Every arm changes exactly one factor of the recipe arm (Table~\ref{tab:arms}); all other settings,
including the seed set, are identical.

\subsection{Lesion-faithful supervision}
Each lesion (LIDC consensus of at least two readers, 50\% voxel agreement) was projected as its excess density over the surrounding lung.
\emph{Lesion-weighted loss} adds $(\mathrm{L1}+\mathrm{MSE})$ on a dilated lesion mask with weight 1. \emph{Counterfactual lesion insertion (CLIA)}
adds, before the physics, 0--3 lesion patches drawn from \emph{training} CTs at random lung positions (including over bone) with random scale
(0.7--1.5) and an amplitude of 0.6--1.6 times that scale; the bone target is untouched, so the network is explicitly taught that a lesion on a rib is not bone.

\subsection{Lesion-absorption audit}
For a model with bone estimate $\hat B$ and a lesion signal $n$ in composite-image units, the absorption is
$A=\langle \hat B(x+\alpha n)-\hat B(x),\,\alpha n\rangle/\langle\alpha n,\alpha n\rangle$ over the lesion support;
$A=0$ means the lesion stays in the soft image and $A=1$ that it was removed with the bone. Because $A=0$ is trivially reached by doing nothing,
it is always reported next to the suppression gain. On synthetic test CTs the lesions are drawn from the test CTs only (never from training CTs) and the truth is exact.
On real radiographs (SZCH test set) lesions from held-out CTs are added at random positions in the lung fields (lung masks from a published chest-X-ray
segmentation network); the bone overlap of a site is measured on the real dual-energy bone signal, never on a model output.

\subsection{Evaluation on real images}
\paragraph{Paired dual-energy data.} The SZCH-X-Rays test set has 75 radiographs with dual-energy-subtraction soft-tissue references.
Tone differs between a radiograph and its DES image, so metrics are reported as is (``raw'') and after a per-image least-squares gain and offset of the estimate
onto the reference (``calibrated'', primary); calibrated PSNR and SSIM are also reported inside the region where DES removed the most signal
(top 10\% of calibrated radiograph minus DES) and outside it.
\paragraph{Rib benchmark.} On VinDr-RibCXR (245 images with per-rib masks; evaluated locally, never used for training) we report the relative reduction of Sobel
gradient energy on the rib-boundary band inside the lungs ($R_\mathrm{rib}$), on lung pixels away from ribs ($R_\mathrm{non}$), their difference (selectivity)
and the ROC-AUC of the gradient-reduction map for separating the two sets of pixels.
\paragraph{Controls and baselines.} Identity, CLAHE, a homomorphic filter and the publicly released CT-supervised weights of the reference method
(research licence) are scored with the same code.

\subsection{Statistics}
All comparisons are paired. Confidence intervals are percentile bootstrap intervals (10\,000 resamples) over patients when several units come from one patient;
tests are Wilcoxon signed-rank tests with Holm correction within a hypothesis family; non-inferiority uses the lower bound of a one-sided 95\% interval against a
margin fixed in the analysis plan; AUCs are compared with the paired DeLong test. Seeds are averaged within unit before testing and their spread is reported.
The analysis plan, the split and the code commit were frozen before any locked evaluation set was scored (\path{freeze.json}).


\section{Details that the paper omits}
\label{sec:dp-details}
The statements in this section are read from the code of the v1 data pipeline, whose files are named in typewriter type. A CT is segmented and projected once into areal-density maps (\path{s4/ct.py},
\path{s4/seg.py}, \path{s4/render.py}); the maps of all usable CTs are loaded onto the GPU as one bank (\path{s4/bank.py}); and every
training sample is composed from the bank at run time (\path{s4/xray.py}, \path{s4/compose.py}). The real radiographs of the downstream tasks
are prepared by a separate script, described first.

\subsection{Data preparation per set}
\label{ssec:dp-prep}
The downstream sets are converted once into fixed arrays by \path{s4/prep.py}, with one loader per set. SZCH-X-Rays and VinDr-RibCXR do not
pass through it; they are read directly by their evaluation scripts (\path{s4/eval_real.py}, \path{s4/eval_ribcxr.py}).

\paragraph{Reading and resizing.} Every image is read as one grey-level array in floating point (MetaImage through SimpleITK, PNG and JPEG
through OpenCV at their full bit depth, colour converted to grey), area-resized to $512\times512$~px without keeping its aspect ratio, as the
renders are (Section~\ref{ssec:dp-geom}), and min--max normalised to $[0,1]$. This is the suppressor input $x$. The variants used downstream---the
original, the suppressed image $x-\hat b$ of the frozen ensemble and optionally the bone estimate and the suppressed image of the released model of
the reference method \citep{angaitkar2026anatomy}, run at its native 1024~px---are clipped to $[0,1]$, area-resized to the target size $s$ and
stored as 8-bit arrays of shape $N\times s\times s$: $s=256$~px for the classification sets (NIH ChestX-ray14, BDCXR-3257) and $s=512$~px for the
detection and nodule sets (TBX11K, VinBigData, Node21, JSRT), whose boxes are rescaled with separate horizontal and vertical factors. For the
detection sets the script also stores, per box, the fraction covered by strong predicted bone (above 25\% of the 99th percentile of the bone map)
and a contrast-to-noise ratio against a surrounding ring of half the box size, for the original, the suppressed and a CLAHE image
\citep{zuiderveld1994clahe}.

\paragraph{Hash splits.} All hash splits in this work, including the LIDC patient split (\path{tools/make_split.py}), map a key $k$ to
\begin{equation}
  h(k) = \mathrm{int}_{16}\big(\mathrm{SHA}_8(\texttt{20261026:}k)\big) \bmod 100,
  \label{eq:dp-hash}
\end{equation}
where $\mathrm{SHA}_8(\cdot)$ is the first 8 hexadecimal digits of the SHA-256 digest of its argument,
and assign it to training if $h<70$, to validation if $70\le h<80$ and to the test set otherwise. The key is the patient identifier where a set has
one and the image identifier otherwise, so the assignment of an image depends on its own key only and survives changes of the file list.

\paragraph{Per set.} \emph{NIH ChestX-ray14} \citep{wang2017chestxray}: 14 multi-hot finding labels; the official test list is the test set, and
the official training and validation list is split by the hash of the patient identifier into validation ($70\le h<80$) and training (all other
values). \emph{TBX11K} \citep{liu2020tbx11k}: boxes from the COCO-format files; the official training images are split by the hash of the file name
into training and validation, and the official validation images form the test set. \emph{VinBigData} \citep{vinbigdata2020kaggle}: images from a
public 1024~px copy, boxes rescaled from the original image size; for each image and finding (``no finding'' is ignored) a box joins the first
cluster whose mean box has an intersection over union above 0.4 with it and that holds no box of the same radiologist. Clusters of at least two
radiologists become one mean box, single-reader clusters are dropped, an image whose findings are all single-reader is excluded, and an image
without findings stays as a normal image; the split is by the hash of the image identifier. \emph{Node21} \citep{sogancioglu2024node21}: the
processed 1024~px images (Zenodo record 5548363), read from a private copy or fetched from the archive by HTTP range requests, with the nodule
boxes halved to 512~px and a hash split by image name; the loader stops if fewer than 99\% of the listed images are present. \emph{JSRT}
\citep{shiraishi2000jsrt}: evaluation only, so every image is a test image; each benign or malignant nodule becomes a square box with the nodule
diameter as side, centred on the marked position (0.175~mm pixels in the 2048~px original, divided by four). \emph{BDCXR-3257}
\citep{durjoy2026bdcxr}: the loader checks the three class counts; the five folds are stratified by class and grouped by near-duplicate group
(Section~\ref{sec:data}), and in each fold whole groups of the training part move to validation until it holds at least 10\% of the training-part images of each
class.

\paragraph{Unreadable images.} An image that cannot be decoded is logged and skipped, never fatal: its array rows stay zero and its name is
listed as unreadable in \path{status.json}. Images that the annotations list but that are missing on disk are logged and left out (TBX11K,
Node21).

\subsection{Segmentation recipe}
\label{ssec:dp-seg}
A CT is loaded by \path{s4/ct.py} as a volume of Hounsfield units indexed $[z,y,x]$, with $z=0$ the most superior slice, clipped to
$[-1024,3071]$~HU; the voxel spacing $(\Delta z,\Delta y,\Delta x)$ travels with it, and every rule below that is stated in millimetres is
converted to voxels with it. \path{s4/seg.py} then builds the body, bone and lung masks in this order, because each mask uses the earlier ones.
The three projected components, bone, lung and the rest of the body, are disjoint.

\paragraph{Body.} (1)~Threshold at $\mathrm{HU}>-300$; (2)~fill the holes of every axial slice; (3)~open with a $1\times9\times9$ voxel
structuring element (in-plane only), which cuts thin connections such as the couch off the body; (4)~keep the largest connected component. The
order of steps 2 and 3 matters: filling first encloses the lungs on every slice with an intact chest wall, whereas an opening first could break a
thin chest wall, and the filling would then leave the lungs out (a failure recorded in the code comment of \texttt{body\_mask}).

\paragraph{Bone.} Bone follows the iterative relaxation of \citet[Algorithm~1]{angaitkar2026anatomy}. Let $R=\{\mathrm{HU}>300\}$ be the
reference mask and $f_M(v)$ the fraction of a mask $M$ inside a cube of side 10~mm centred on voxel $v$. Starting from $M_0=R$, three steps admit
lower thresholds only next to bone already found:
\begin{equation}
\begin{gathered}
M_k = \big(\{\mathrm{HU}>\tau_k\}\cap\{v: f_{M_{k-1}}(v)\ge\theta\}\big)\cup R,\\
(\tau_1,\tau_2,\tau_3)=(200,150,100)~\mathrm{HU},
  \quad \theta=0.05 .
\end{gathered}
\label{eq:dp-bone}
\end{equation}
A voxel of 120~HU is bone only if at least 5\% of its 10~mm neighbourhood was bone after the second step, so the mask grows from the cortex into
trabecular bone without admitting isolated soft-tissue voxels of the same attenuation. To keep this affordable, $f_M$ is computed on a grid
coarsened by $\max(1,\mathrm{round}(3~\mathrm{mm}/\Delta))$ along each axis and copied back; voxels of an incomplete last block receive the
fraction zero and can enter only through $R$. The mask is then closed with a ball of radius two voxels, which fills the marrow inside the
cortical shell; the closing runs on the bounding box of the mask, edge-padded by 3 voxels because a binary closing in SciPy erodes structures
that touch the array border. Components smaller than $\kappa=1\%$ of the largest one are removed (the code calls them speckle islands), and the
renderer intersects the result with the body mask.

\paragraph{Lung.} Air inside the body, $\{\mathrm{HU}<-320\}\cap M_{\mathrm{body}}$, is split into connected components, and every component with
at least 20\% of the voxels of the largest one is kept; this keeps both lungs, connected or not, and drops small gas pockets. On every slice that
contains lung the mask is closed with a disc of radius 5 pixels and hole-filled, so that vessels and nodules enclosed by aerated lung belong to the
lung component, and bone voxels are removed. The trachea and main bronchi are connected air and therefore part of the lung mask, which matters
for the frame (Section~\ref{ssec:dp-geom}).

\paragraph{Soft tissue.} The non-lung soft tissue is $M_{\mathrm{soft}}=M_{\mathrm{body}}\setminus(M_{\mathrm{lung}}\cup M_{\mathrm{bone}})$.
If the bone voxels were simply removed from it, the soft-tissue projection would carry a negative imprint of the skeleton, as if the bones were
air. Its density inside the bone mask is therefore filled from the surrounding soft tissue by normalised convolution,
\begin{equation}
  \rho_{\mathrm{fill}}(v) = \frac{\big(G_\sigma*(\rho\,\mathbf{1}_{M_{\mathrm{soft}}})\big)(v)}{\big(G_\sigma*\mathbf{1}_{M_{\mathrm{soft}}}\big)(v)},
  \label{eq:dp-fill}
\end{equation}
with a Gaussian $G_\sigma$ of $\sigma=4$~mm where the denominator exceeds 0.01, then $\sigma=12$ and 36~mm for voxels still unfilled deep inside
large bones, and the density of water for any voxel left. The lung component keeps the measured density of its voxels, vessels and nodules
included. The density $\rho=\mathrm{clip}(1+\mathrm{HU}/1000,0,3)$~g\,cm$^{-3}$ treats every voxel as water of a different density; since HU
measure attenuation relative to water at the effective energy of the CT beam, this mixes density and composition for bone, whose composition
enters only through the cortical-bone attenuation table (Section~\ref{ssec:dp-atten}). The simplification relative to the reference method is
recorded in \path{DEVIATIONS.md} (2026-10-26).

\paragraph{Sources of the rules.} The bone recipe follows the published method \citep[Algorithm~1]{angaitkar2026anatomy} in its structure:
a reference mask above 300~HU, a threshold ladder of 200, 150 and 100~HU admitted by support voting with $\theta=0.05$, and pruning of
components smaller than 1\,\% of the largest. It differs in three details: support is counted in a 10~mm cube instead of a window of $15^3$
voxels, the first pass starts from the reference mask instead of from all voxels, and the reference method's closings at every step (balls of
radius 1 and 2) are replaced by one radius-2 closing after the ladder. The body, lung and soft-tissue steps are choices of this work: the
reference method removes the bed and clips intensities, segments the lungs with a pretrained model, and fills the bone voxels of the soft-tissue
volume by diffusion inpainting; this work uses the morphology listed above (fill-then-open order, $1\times9\times9$ opening, 20\,\% rule,
radius-5 closing), a coarse grid and a normalised-convolution fill.

\paragraph{Failure modes.} The code handles three: the fill-then-open order of the body mask, the border erosion of SciPy's closing, and
trachea-only slices at the top of the scan, which the frame does not count as lung (Section~\ref{ssec:dp-geom}). Others follow from the thresholds
and are not detected. Any voxel above 300~HU becomes bone whatever its anatomy, so calcified nodules or lymph nodes, densely contrast-filled vessels
and metal would enter the bone target; lung disease denser than $-320$~HU belongs to the lung component only where the slice-wise closing encloses
it; and a couch that the opening does not detach stays in the body. Gross failures are screened by the usable-CT rules of
Section~\ref{sec:methods}, which \path{tools/make_split.py} applies to the metadata of every render and records as named reasons in the split
file; the lungs count as cut by the scan when their extent (Section~\ref{ssec:dp-geom}) starts in the first two or ends in the last two slices.
How often the threshold-based failures occurred was not measured: no step of the pipeline counts calcified nodules, contrast or metal in the
bone target, so their effect on the v1 models is unknown. The v2 models are trained on the Qure.ai projections and do not depend on this
segmentation.

\subsection{Rendering geometry}
\label{ssec:dp-geom}
\paragraph{Orientation and sources.} The DICOM loader accepts only axial series: the orientation of the first slice must equal $(1,0,0,0,1,0)$ to
within 0.02 in every direction cosine. It sorts the slices by position, puts the head first and skips series with fewer than 80 CT slices,
duplicate positions, a spacing that deviates anywhere from the median by more than 25\%, slices of unequal size or an unreadable file; a skipped
CT is logged and never stops the run. With this orientation the row index $y$ runs from anterior to posterior and the column index $x$ from the
patient's right to the patient's left, so a projection along $y$ shows the head at the top and the patient's right on the viewer's left, the
display convention of a posteroanterior radiograph.

Each series is read from the mounted DICOM copy (606 of the 1\,018 series, \path{DEVIATIONS.md}, 2026-10-26), otherwise from the LUNA16
MetaImage copy \citep{setio2017luna16}, otherwise downloaded from TCIA (\path{kernels/render/render_shard.py}); the source is recorded with
each render. The MetaImage loader accepts uncompressed 16-bit volumes only and locates a nodule slice from its $z$ position in millimetres, whereas
the DICOM loader matches the SOP instance identifier of the annotation. The agreement of the two copies was checked by
\path{kernels/orient_check/check.py} on 12 series spread over those present in both: shapes, spacings, the mean absolute HU difference for all
eight combinations of axis flips and the first slice of every consensus nodule located either way. The outcome is given in
Section~\ref{sec:methods}.

\paragraph{Projection and the AP view.} Each component is projected by a parallel-beam sum along the anteroposterior axis,
\begin{equation}
  a_c[z,x]=\frac{\Delta y}{10}\sum_y\rho_c[z,y,x]\quad(\mathrm{g\,cm^{-2}},\ \Delta y\ \text{in mm}),
  \label{eq:dp-project}
\end{equation}
which is the same for a beam entering from the front (AP) or the back (PA); the raw images of the two views differ by a horizontal mirror. The
renderer stores the PA convention, and the composer draws a horizontal mirror with probability 0.5 inside the affine transform that it shares
between all maps (\texttt{random\_affine} in \path{s4/compose.py}), so that both handednesses occur in training. Divergent-beam magnification,
which does differ between AP and PA, is not modelled.

\paragraph{Frame.} The frame (\texttt{frame\_box} in \path{s4/render.py}) is a box in the $[z,x]$ plane of the projection:
\begin{itemize}
\item Lung extent $L$: from the first to the last slice whose lung cross-section exceeds 8~cm$^2$. Trachea-only slices must not count, or every
scan that starts in the neck would look cut at the apex and the frame would be centred too high.
\item Width: $W=1.08$ times the left--right extent of the body mask over all slices, centred on the body.
\item Height: $H=\mathrm{clip}(L+60~\mathrm{mm},\,0.75\,W,\,1.2\,W)$, centred on the lungs.
\item Placement: if the scan is at least $H$ long, the frame is slid along $z$ until it lies inside the scan, so that no row is padded; otherwise
it stays centred and the rows beyond the scan are padded.
\end{itemize}
The code comment calls $0.75\le H/W\le1.2$ the aspect range of real radiographs and notes that short LIDC scans cannot reach the upper end.
No source for this range was found; it is a design choice of this work.
The box is area-resized to $512\times512$ pixels, so one pixel covers $W/512$~mm horizontally and $H/512$~mm vertically: the frame is squashed to
a square exactly as real radiographs are squashed by \path{s4/prep.py}, and a render and a radiograph of the same aspect ratio carry the same
distortion. The two-dimensional lung and body masks and the lesion signals are framed in the same way with zero padding; the background of a lesion
signal is the median density of a three-voxel ring around the nodule inside the lung and outside bone.

\paragraph{Padded rows.} Where the frame extends beyond the scan, the renderer repeats the first or last projected row (edge padding) and records
the number of padded frame rows at the top and bottom, $\mathrm{round}(512\,p/(z_1-z_0))$ for $p$ padded slices of a frame spanning $z_1-z_0$
slices. Edge padding turns the last projected row into vertical streaks. On 2026-10-27 two rules changed together (\path{DEVIATIONS.md}): of the
renders, 38\% needed more than the original limit of 25 padded rows because their scans are shorter than $0.75\,W$, so the usability limit was
raised to 50 rows in total (10\% of the frame); and the padded rows are set to air, zero areal density in all three components, when the renders
are loaded into the bank (\path{s4/bank.py}), instead of keeping the streaks. The renders on disk are unchanged. A row of air has no bone, so its
bone target stays zero under every tone curve (Section~\ref{ssec:dp-gap}).

\subsection{Attenuation and spectra}
\label{ssec:dp-atten}
\paragraph{Materials and tables.} The bone, non-lung soft-tissue and lung components are given the cortical bone, soft tissue and lung tissue of
ICRU Report~44 \citep{icru1989report44}. Cortical bone differs from the other two by its calcium and phosphorus content, which raises its
photoelectric absorption at low energy; soft tissue and lung tissue have nearly the same composition, and their tabulated mass attenuation
coefficients differ by less than 2\% over the energy grid used here, so the lung and soft-tissue components differ almost only through density.
Table~\ref{tab:icru44} lists the compositions as tabulated by NIST \citep{hubbell2004xraymass}. The renderer itself does not use the nominal
densities: it takes the density of every voxel from the CT (Section~\ref{ssec:dp-atten}), and the composition enters only through the
attenuation table of each material.
\begin{table}[htbp]
\centering\small
\caption{ICRU-44 compositions used through the NIST attenuation tables: density and elemental weight fractions (NIST SRD 126, Table~2).}
\label{tab:icru44}
\begin{adjustbox}{max width=\linewidth}\begin{tabular}{@{}lccccccccccccc@{}}
\toprule
Material & $\rho$ (g/cm$^3$) & H & C & N & O & Na & Mg & P & S & Cl & K & Ca \\
\midrule
Cortical bone & 1.920 & 0.034 & 0.155 & 0.042 & 0.435 & 0.001 & 0.002 & 0.103 & 0.003 & -- & -- & 0.225 \\
Soft tissue & 1.060 & 0.102 & 0.143 & 0.034 & 0.708 & 0.002 & -- & 0.003 & 0.003 & 0.002 & 0.003 & -- \\
Lung tissue & 1.050 & 0.103 & 0.105 & 0.031 & 0.749 & 0.002 & -- & 0.002 & 0.003 & 0.003 & 0.002 & -- \\
\bottomrule
\end{tabular}\end{adjustbox}
\end{table}
The coefficients $(\mu/\rho)_c(E)$ are the NIST tabulations for these materials (Table~4 of the NIST X-ray mass attenuation coefficient pages,
which rest on the XCOM cross sections \citep{berger2010xcom}).
These pages are the NIST X-ray mass attenuation coefficient tables \citep{hubbell2004xraymass}, not the XCOM program.
\path{tools/build_xray_tables.py} downloads the three pages, keeps them in \path{s4/data/nist_raw/} for provenance, keeps the rows between 8
and 200~keV, separates the duplicate energies at absorption edges by $10^{-6}$~keV and interpolates $\ln(\mu/\rho)$ linearly in $\ln E$ onto a
1~keV grid from 10 to 150~keV. The result is committed as \path{s4/data/xray_tables.json}, so no run depends on the NIST server. On this grid
the ratio of the bone to the soft-tissue coefficient is about 3.5 at 30~keV, 1.5 at 60~keV and 1.1 at 100~keV, which is why bone contrast falls
with tube voltage.

\paragraph{Spectra.} Tube spectra come from SpekPy \citep{poludniowski2021spekpy} for a tungsten anode at $12^\circ$, with 1~keV bins, the
\texttt{kqp} physics option and 2.5~mm of aluminium filtration, at 60 to 140~kVp in 10~kV steps. Each spectrum is interpolated onto the
10--150~keV grid and normalised to unit sum; the SpekPy version is recorded in the provenance field of the table file. Photons below 10~keV,
which the aluminium filter makes negligible, are not represented.

\paragraph{Detector and lookup tables.} The detector is energy-integrating, so each photon is weighted by its energy. With the normalised weight
$p(E)=\phi(E)\,E/\sum_{E'}\phi(E')\,E'$, the pixel value of component $c$ at areal density $a$ is
\begin{equation}
  P_c(a) = -\ln\sum_E p(E)\,\exp\!\big(-(\mu/\rho)_c(E)\,a\big),
  \label{eq:dp-pixel}
\end{equation}
as in Section~\ref{sec:methods}. For each voltage and material, $P_c$ is tabulated at 8\,192 points from 0 to 80~g\,cm$^{-2}$ and interpolated
linearly on the GPU, with larger areal densities clamped to the end of the table. The renders store areal densities, so one render serves every
voltage. The composer holds tables for 70 to 140~kVp; training draws one of 80, 90, \ldots, 130~kVp per sample with equal probability
(\path{s4/train.py}), and validation uses 100~kVp.

\subsection{Additivity gap and tone curves}
\label{ssec:dp-gap}
\paragraph{Why the components are not additive.} For the whole body, with areal densities $\mathbf a=(a_c)_c$ along a ray, the pixel value is
\begin{equation}
  P(\mathbf a) = -\ln\sum_E p(E)\,\exp\!\Big(-\sum_c(\mu/\rho)_c(E)\,a_c\Big),
  \label{eq:dp-whole}
\end{equation}
while the composer uses the separate $P_c(a_c)$ of Equation~\eqref{eq:dp-pixel}. For a monochromatic beam at $E_0$ both are linear and
$P(\mathbf a)=\sum_c(\mu/\rho)_c(E_0)\,a_c=\sum_cP_c(a_c)$: the components add exactly. For a polychromatic beam they do not. Writing
$p_a(E)\propto p(E)\exp(-(\mu/\rho)_c(E)\,a)$ for the spectrum after $a$~g\,cm$^{-2}$ of material $c$,
\begin{equation}
\begin{gathered}
P_c'(a)=\mathrm{E}_{p_a}\big[(\mu/\rho)_c\big]>0,\\
P_c''(a)=-\mathrm{Var}_{p_a}\big[(\mu/\rho)_c\big]\le0 ,
\end{gathered}
\label{eq:dp-concave}
\end{equation}
so each further gram per square centimetre attenuates less than the previous one, because the low-energy photons, which are attenuated most, go
first: beam hardening. The same mechanism couples the components: inside the body, the beam that reaches the bone has already been hardened by
the surrounding soft tissue, and the bone attenuates it less than it would attenuate the unhardened beam.

\emph{Claim.} If every $(\mu/\rho)_c$ is non-increasing in $E$ on the support of $p$, then $P(\mathbf a)\le\sum_cP_c(a_c)$: the sum of separately
projected components overestimates the attenuation of the whole body.

\emph{Proof.} The functions $f_c(E)=\exp(-(\mu/\rho)_c(E)\,a_c)$ are non-negative and non-decreasing. For two such functions $f,h$, Chebyshev's sum
inequality with weights $p$ gives $\sum_Ep\,f\,h\ge(\sum_Ep\,f)(\sum_Ep\,h)$; a product of non-negative non-decreasing functions is again of this
kind, so induction over the components gives $\sum_Ep\prod_cf_c\ge\prod_c\sum_Ep\,f_c$, and applying the decreasing function $-\ln$ to both sides
yields the claim.

The condition holds for the three tables of Section~\ref{ssec:dp-atten}, whose coefficients decrease monotonically on the 10--150~keV grid,
above the absorption edges of the constituent elements. The size of the gap for one representative ray is given in Section~\ref{sec:methods}, and
\path{tests/test_renderer.py} checks it against a fixed gate.
Figure~\ref{fig:additivity_gap} shows the gap over the range of the renders: it reaches about 8\,\% of the joint projection at 80~kVp, falls
with the tube voltage, and grows with the bone areal density.
\begin{figure}[htbp]
\centering
\safefig{figures/additivity_gap.pdf}
\caption{Beam-hardening gap of separately projected components, $(P_{\mathrm{bone}}+P_{\mathrm{soft}}-P)/P$, against the soft-tissue areal
density for three bone areal densities and three tube voltages (attenuation tables and spectra of the renderer; \texttt{tools/paper\_extras.py}).}
\label{fig:additivity_gap}
\end{figure}

\paragraph{What the composer does with the gap.} The composer never forms $P(\mathbf a)$. With per-sample quantiles $q_\alpha$ it adds weighted
component images,
\begin{equation}
\begin{gathered}
S=\frac{1}{1.5}\,\frac{s-q_{0.01}(s)}{q_{0.995}(s)-q_{0.01}(s)},\quad s=P_{\mathrm{nl}}+g\,P_{\mathrm{lung}},\\
B=\frac{1}{1.5}\,\frac{w\,P_{\mathrm{bone}}}{q_{0.99}(P_{\mathrm{bone}})},\\
X=S+B,
\end{gathered}
\label{eq:dp-compose}
\end{equation}
with $g\sim U(1,3)$, $w\sim U(0.2,0.5)$ and $1.5=1+\max w$ keeping $X$ within about $[0,1]$. The input is thus exactly additive in the logarithmic
domain by construction, and the bone target is exactly its bone term. The gap separates the composed input from a physically exact radiograph of
the same CT; it is not an error of the target. The design trades physical exactness of the input, which the random weights and the quantile
normalisation loosen further, for exactness of the target.

\paragraph{Tone curves.} A clinical radiograph is not the log-attenuation image itself: detector response, vendor processing and display apply a
monotone grey-level transform. The composer draws one transform $\gamma$ per sample (\texttt{apply\_tone} in \path{s4/compose.py}), piecewise
linear between knots $(\xi_j,\eta_j)$ and constant outside them:
\begin{itemize}
\item \emph{Gamma curve}: $\xi_j=1.5\,j/64$ for $j=0,\ldots,64$ and $\eta_j=\xi_j^{\,e}$ with $e\sim U(0.7,1.4)$.
\item \emph{Histogram matching}: $\xi_j$ are the 257 quantiles of $X$ at probabilities $j/256$ and $\eta_j$ the 257 stored quantiles of a real
radiograph drawn from the reference set of Section~\ref{sec:methods}, both plus a ramp from 0 to $10^{-5}$ that makes the knots strictly increasing.
\end{itemize}
In both cases $\eta_j$ is non-decreasing, so $\gamma$ is non-decreasing on the whole real line; the curve fitted to $X$ is also applied to $S$.
Training uses histogram matching unless an arm says otherwise; validation uses the gamma curve without detector effects or lesion insertion, so
that all arms are scored on the same images.

\paragraph{Exactness under the tone curve.} The composer sets
\begin{equation}
\begin{gathered}
X'=\gamma(X),\\
S'=\gamma(S),\\
B'=X'-S' .
\end{gathered}
\label{eq:dp-tone}
\end{equation}
\emph{Claim.} For every non-decreasing $\gamma$: (i) $X'=S'+B'$ exactly; (ii) $B'\ge0$; (iii) $B'=0$ wherever $B=0$; (iv) if $\gamma$ is
continuous and piecewise linear with largest slope $\lambda$, then wherever $B>0$, $B'=\bar\gamma'B$ with
$\bar\gamma'=B^{-1}\int_S^{S+B}\gamma'(t)\,\mathrm dt\in[0,\lambda]$.

\emph{Proof.} (i) holds by the definition of $B'$, for any $\gamma$. (ii) $P_{\mathrm{bone}}\ge0$ because $P_c(0)=0$ and $P_c'>0$
(Equation~\eqref{eq:dp-concave}); $w>0$ and the quantile in the denominator is bounded below by $10^{-6}$ in the code, so $B\ge0$, $X\ge S$ in
every pixel and $\gamma(X)\ge\gamma(S)$. (iii) $B=0$ gives $X=S$, hence $B'=0$. (iv) A continuous piecewise-linear $\gamma$ is absolutely
continuous, so $\gamma(S+B)-\gamma(S)=\int_S^{S+B}\gamma'(t)\,\mathrm dt$ with $0\le\gamma'\le\lambda$ almost everywhere.

Two consequences matter. First, $B'$ is in general not $\gamma(B)$ but the bone signal times the mean slope of the curve over the grey levels the
bone spans in that pixel, so the same bone appears stronger where the curve is steep; it is exactly the quantity whose subtraction leaves the
tone-mapped soft tissue, $X'-B'=\gamma(S)$, which $\gamma(B)$ would not be for a non-linear curve. Second, by (iii) the target is zero wherever the
bank holds no bone, the padded rows included.

\paragraph{Detector effects.} With detector effects on, the default for training, three effects follow the tone curve. Let $X'_b=G_\sigma*X'$
and $B'_b=G_\sigma*B'$, with $G_\sigma$ a separable Gaussian (replicate padding) whose $\sigma\sim U(0,1.2)$~px is drawn per sample with
probability 0.5 and applied when $\sigma>0.05$~px. Then
\begin{equation}
\begin{gathered}
X'_{\mathrm{det}} = X'_b + u_1\big(0.004+0.02\sqrt{\max(X'_b,0)}\big)\,\varepsilon + 0.08\,u_2\,\big(G_{20}*X'_b\big),\\
B'_{\mathrm{det}} = B'_b ,
\end{gathered}
\label{eq:dp-detector}
\end{equation}
with $\varepsilon\sim\mathcal N(0,1)$ per pixel and $u_1,u_2\sim U(0,1)$ per sample: noise whose standard deviation grows with the square root of
the intensity above a floor, and a low-frequency copy of the image (Gaussian of 20~px) with up to 8\% of its amplitude as scatter. Only the blur
touches the bone target, and because it is linear and identical on $X'$ and $B'$, $X'_b-B'_b=G_\sigma*S'$ and the decomposition stays exact.
Noise and scatter are added to the input only, so the image the network should leave after subtracting its bone estimate contains them: it is
taught to keep noise and scatter in the suppressed image. The loss uses the bone target only; the returned soft-tissue image is $S'$ before the
detector effects and serves validation, which runs without them. The code names no source for the noise and scatter constants; they are
augmentations of plausible form, not a calibrated detector model.

\subsection{Tests}
\label{ssec:dp-tests}
The pipeline has 32 unit tests in eight files under \path{tests/}. None needs data, a GPU or a network connection: every test builds its own
input (a CT phantom, a minimal LIDC annotation file, a MetaImage volume, a small fake projection pack or random scores), and the only repository
file they read is the committed \path{s4/data/xray_tables.json}. They run on the CPU from the repository root with
\begin{verbatim}
python -m pytest tests -q
\end{verbatim}
and need NumPy, SciPy, scikit-image, OpenCV, pydicom, PyTorch, timm, scikit-learn and pytest; \path{pytest.ini} sets the test path and writes
temporary files to \path{outputs/pytest_tmp}.

\paragraph{\texttt{test\_renderer.py} (7 tests).} A phantom of $120\times160\times200$ voxels ($2.5\times1.5\times1.5$~mm) with an elliptic body of
40~HU, two elliptic lungs of $-850$~HU, a spine cylinder of 600~HU and four rib rings of 700~HU gives closed-form answers. The tests check that
(1)~the bone mask holds more than 99\% of the voxels above 300~HU and none below $-500$~HU; (2)~the lung volume is within 15\% of the analytic
one; (3)~the projection of a cylinder of 1.6~g\,cm$^{-3}$ equals chord length times density to 8\%; (4)~component pixel values add exactly (to
$10^{-3}$) for a 60~keV monochromatic beam, and at 100~kVp the relative gap for 20~g\,cm$^{-2}$ of soft tissue and 3~g\,cm$^{-2}$ of bone stays
below a gate of 0.15 fixed in advance (the test prints the value; the paper reports it); (5)~every lookup table starts at zero and is strictly
increasing; (6)~for a nodule of 30~HU in lung of $-850$~HU the background density is 0.15~g\,cm$^{-3}$ to 0.03 and the peak excess areal density
is within 12\% of its closed form; (7)~a full render gives three finite, non-negative $512\times512$ maps with bone and uncut lungs, and a volume of
air raises the skip exception.

\paragraph{\texttt{test\_compose.py} (6 tests).} (1)~The GPU lookup agrees with the NumPy evaluation of Equation~\eqref{eq:dp-pixel} to
$2\times10^{-4}$; (2)~the affine transform without rotation, scaling, shift and mirror is the identity; (3)~lesion insertion never changes the
bone or non-lung soft-tissue map and adds lesion signal and mask; (4)~$X'=S'+B'$ holds to $10^{-6}$ with $B'\ge-10^{-6}$ at 80, 100, 120 and
140~kVp; (5)~histogram matching gives the composed image the median of the reference to 0.02 and keeps $B'$ non-negative; (6)~the suppressor runs
forward and backward at 256~px with finite loss and gradients, and the SSIM of an image with itself is 1.

\paragraph{\texttt{test\_io.py} (5) and \texttt{test\_qp.py} (2).} On LIDC annotation files written by the test, the series identifier is read,
two readers form one consensus nodule with the right reader count, malignancy and slice, one reader is not a consensus, and of three readers the
distant one is left out; a synthetic MetaImage volume is put head first, located by $z$ position, and rejected if shorter than 80 slices. On a fake
pack in the format of the v2 source (Section~\ref{sec:method_v2}), $X'=S'+B'$ holds to $10^{-5}$ with a non-negative bone image, lesion insertion
changes the input while the bone target keeps more than 90\% of its sum, and with detector effects on, $X'-B'-S'$ is no longer zero, because noise
and scatter reach the input only.

\paragraph{Statistics tests (12).} The other four files check the evaluation code on constructed cases: the paired DeLong test
\citep{delong1988auc} against the AUCs of scikit-learn and its rejection rate under the null (\path{test_delong.py}); FROC, average precision
and one-to-one box matching on known boxes (\path{test_froc.py}); the Holm correction \citep{holm1979simple}, clustered bootstrap intervals,
the Wilcoxon test \citep{wilcoxon1945individual} and the non-inferiority rule (\path{test_stats.py}); and conformal coverage and sparsification
error of the uncertainty estimates (\path{test_uncertainty.py}).

The tests cover the steps whose errors would silently change the supervision. Beyond the single phantom render they do not check the frame placement
and padded-row counts, the per-set loaders of \path{s4/prep.py} or the GPU bank; these are checked only through the counts and logs that their
runs print.
Two parts of the pipeline have no unit test: the frame box (padding counts on a short scan) and the VinBigData consensus rule.
